% Generated draft: requires mathematical and source review.
\documentclass[11pt]{article}
\usepackage[margin=25mm]{geometry}
\usepackage{fontspec}
\setmainfont{Latin Modern Roman}
\usepackage{amsmath,amssymb,mathtools,unicode-math}
\setmathfont{Latin Modern Math}
\usepackage{enumitem,needspace,xurl,hyperref}
\setlength{\parindent}{0pt}
\setlength{\parskip}{5pt}
\setlength{\emergencystretch}{3em}
\newcommand{\N}{\mathbb{N}}
\newcommand{\Z}{\mathbb{Z}}
\newcommand{\Q}{\mathbb{Q}}
\newcommand{\R}{\mathbb{R}}
\newcommand{\C}{\mathbb{C}}
\newcommand{\F}{\mathbb{F}}
\newcommand{\A}{\mathbb{A}}
\newcommand{\PP}{\mathbb{P}}
\newcommand{\E}{\mathbb{E}}
\newcommand{\eps}{\varepsilon}
\newcommand{\dd}{\,\mathrm{d}}
\begin{document}
\section*{UnsolvedMath ID: 700085}
\textbf{Source problem code:} catalogue entry
\subsection{Definitions and mathematical statement}
Let \(G\) be a group. Write \(\mathcal C_{\mathrm{gen}}(G)\) for the set called the \emph{generalized central elements} of \(G\) in the supplied source, and write \(\mathcal N_{\mathrm{nil}}(G)\) for its \emph{nil-kernel}. These symbols are placeholders for the source terminology: the supplied record defines neither of these notions nor the term \(NR\)-group. Consequently, the formulation below requires those definitions before it can be a self-contained mathematical statement.

The supplied record states the forward implication
\[
G\text{ is an }NR\text{-group}
\quad\Longrightarrow\quad
\mathcal C_{\mathrm{gen}}(G)=\mathcal N_{\mathrm{nil}}(G).
\]
The question, attributed to V. G. Vilyatser, asks whether the converse holds: for every group \(G\), does
\[
\mathcal C_{\mathrm{gen}}(G)=\mathcal N_{\mathrm{nil}}(G)
\quad\Longrightarrow\quad
G\text{ is an }NR\text{-group}?
\]
The record specifies no restriction to finite groups or to another class of groups; any restrictions implicit in the undefined terminology remain unresolved.
\subsection{Short English statement}
V. G. Vilyatser asks whether equality of a group's generalized central elements and its nil-kernel forces the group to be an \(NR\)-group. The supplied record states the reverse implication. Definitions of all three specialized notions are missing, so this draft retains the question provisionally and requires mathematical review.
\subsection{Sources}
\begin{itemize}
\item Local UnsolvedMath registry, supplied record with permanent ID \(700085\), exact source problem code \texttt{catalogue entry}, title \emph{Kourovka Notebook Problem 2.6}; upstream revision \texttt{d371ee917ac9647fbb06cd6fc529deec6f3dda82}.
\item Supplied bibliographic attribution: \emph{The Kourovka Notebook (Unsolved Problems in Group Theory)}, problem 2.6, second issue (1967), proposed by V. G. Vilyatser; cited in the record through the twenty-first issue edition (2026), \texttt{arXiv:1401.0300v48}. The supplied collection description names E. I. Khukhro and V. D. Mazurov as editors.
\item The local record labels the problem \emph{open}. This is supplied registry status, not an independently verified literature assessment. The cited source and its terminology have not been independently checked for this draft.
\end{itemize}
\end{document}
